\documentclass[10pt,a4paper]{amsart}
\usepackage[margin=19mm]{geometry}
\usepackage{amsmath,amssymb,mathtools}
\usepackage[scr=rsfs,cal=euler]{mathalfa}
\usepackage{xfrac}
\usepackage[unicode,hidelinks,bookmarks=false]{hyperref}
\newcommand{\ie}{i.e.\ }
\newcommand{\cf}{cf.\ }
\newcommand{\resp}{respectively\ }
\newcommand{\Subsection}{\subsection}
\providecommand{\nobreakdash}{\nobreak\hbox{-}}
\setlength{\emergencystretch}{2em}
\multlinegap0pt
\usepackage{bm}
\usepackage{bbm}
\usepackage{dsfont}
\usepackage{xspace}
\usepackage{cancel}
\usepackage{mathtools}
\usepackage{amsthm}
\makeatletter
\@ifundefined{theorem}{\newtheorem{theorem}{Theorem}}{}
\@ifundefined{proposition}{\newtheorem{proposition}[theorem]{Proposition}}{}
\makeatother
\newcommand{\Dc}{D^{\circ}}
\newcommand{\Ess}{\mathrm{Ess}}
\newcommand{\Gvwt}{G_{v,w}}
\newcommand{\Nvw}{\mathcal{N}_{v,w}}
\newcommand{\SW}{\mathrm{SW}}
\newcommand{\cn}{\nu}
\newcommand{\rk}{r}
\title[Torus Actions on Matrix Schubert and Kazhdan-Lusztig Varieti]{Torus Actions on Matrix Schubert and Kazhdan-Lusztig Varieties, and their Links to Statistical Models}
\author{Elke Neuhaus, Irem Portakal,, Niharika Chakrabarty Paul}
\date{Source: arXiv:2510.00250v2 (CC BY 4.0)}
\begin{document}
\maketitle

\noindent\emph{Source and licence.} Torus Actions on Matrix Schubert and Kazhdan-Lusztig Varieties, and their Links to Statistical Models. Version arXiv:2510.00250v2 (CC BY 4.0), distributed under the Creative Commons Attribution 4.0 International licence, as recorded in the arXiv licence metadata for this exact version and shown on the abstract page https://arxiv.org/abs/2510.00250v2. Full rights notice: licenses/arxiv-2510.00250-CC-BY-4.0.txt.\par\smallskip

\noindent\textit{Editorial scope.} Counted unit: the Proposition whose source label is \texttt{prop: rectangle comp} in the article (source0.tex lines 1680--1686, source label \texttt{prop: rectangle comp}), delivered together with the article's own complete original proof of it (source0.tex lines 1687--1704, one \texttt{proof} environment of 18 lines). No mathematics is written, repaired or re-derived: statement and proof are the article's own text line for line. Required context retained uncounted: none. Every definition, display and label of those lines is retained unchanged. Omitted from this package and not redistributed: 1--1679, 1705--2478. Nothing inside a retained excerpt is omitted or altered: every retained body is delivered exactly as the article's lines are written, so no omission receipt of any kind accompanies this package. Citations: the retained excerpts carry no citation command at all, so no bibliography entry of the article is transcribed and no reference is redistributed. Reference adaptation: none was needed and none was made; every reference inside the delivered unit resolves inside it, so no unresolved reference or citation is produced. Printed slips kept literal: no printed word, symbol, hypothesis or number of the counted statement or of its proof was altered. Layout and class: the article's own class is \texttt{amsart} with packages including geometry, fontenc, amsthm, amsmath, hyperref, hypernat, comment, graphicx, subfigure, latexsym, amssymb, float, epsfig, multirow; the wrapper is the lane's pinned \texttt{amsart} editorial wrapper at 10pt on A4, which restates the theorem-like environments the retained text needs and adds the article's own macro definitions that the retained text needs (\texttt{\textbackslash Dc}, \texttt{\textbackslash Ess}, \texttt{\textbackslash Gvwt}, \texttt{\textbackslash Nvw}, \texttt{\textbackslash SW}, \texttt{\textbackslash cn}, \texttt{\textbackslash rk}), with extra wrapper packages (bm, bbm, dsfont, xspace, cancel, mathtools). The delivered numbering is the unit's own. Assets: no retained span contains a figure environment, a graphics-inclusion command, a PDF-page inclusion, a table or a TikZ picture; no image file is copied into this package, the package declares no asset dependency and no image file is redistributed with it. Licence: version arXiv:2510.00250v2 of this article is distributed under the Creative Commons Attribution 4.0 International licence, as recorded in the arXiv licence metadata for this exact version and shown on the abstract page https://arxiv.org/abs/2510.00250v2. A full rights notice is delivered at licenses/arxiv-2510.00250-CC-BY-4.0.txt. Characters: every retained excerpt is pure ASCII, so the delivered engine, XeLaTeX with its default Latin Modern text font, reports no missing character.
\begin{proposition}\label{prop: rectangle comp}
    Consider $w\in S_n$, such that $\Dc(w)$ consists of a single rectangle of size $l \times k$. Denote by (a,b) the unique element in $\Ess(w)$ and $m := r_w(a,b)$. Also, denote by $p$ the complexity of $\Nvw$.
    \begin{enumerate}
        \item If $\rk_v(a,b) < m$ then $p= \cn(\Gvwt) - |\Dc(w)| =  |P_v|- kl$. 
        \item If $\rk_v(a,b) =m$ then $p = \cn(\Gvwt)$ and $\Nvw$ is toric if and only if $\Gvwt$ is a forest.
    \end{enumerate}
    \end{proposition}

    \begin{proof}
     First, for $c,d \in [n]$, denote by $\rk_m(c,d)$, the maximal subrank $\rk_u(c,d)$ of matrices $u \in S_n$ with $\rk_u(a,b) \leq m$. We claim that for all such $c,d$, it is $\rk_w(c,d) = \rk_m(c,d)$. Notice that $\rk_m(c,d) \leq \min \{ n-c+1,d \}$. Meanwhile, $w$ must be of the following form:
     $$
    w= (n, \ldots, n-m+1, \, n-m-l, \ldots, n-m-l-k+1, \, n-m, \ldots, n-m-l+1, \, n-m-l-k, \ldots, 1).
$$
     If $c<a$ and $d>b$, then $\rk_w(c,d) = \min \{n-c+1,d \}$. If $c \geq a$ and $d \leq b$, then both $\rk_m(c,d)$ and $\rk_w(c,d)$ are given by $\min \{ m, n-c+1, d \}$.
     If wlog $c \geq a$ and $d>b$, then both $\rk_m(c,d)$ and $\rk_w(c,d)$ are given by $\min \{ m+(d-c), n-c+1 \}$.
     Hence, $w$ dominates all permutations $u$ with $\rk_u(a,b) \leq m$ in the Bruhat order. 
     
    \noindent
    Recall that the complexity of $\Nvw$ is given by $\cn(\Gvwt) - |\Dc(w)| + \# (\textrm{unexpected zeros})$ and that the unexpected zeros of $\Sigma_v$ are the $(i,j) \in \Dc(v)$ such that $t_{v(j),i}v \not\leq w$. We study the change of $\rk_v(a,b)$ under multiplication of $v$ with $t_{v(j),i}$. In the permutation matrix of $v$, left multiplication by $t_{v(j),i}$ swaps rows $v(j)$ and $i$. More precisely, it interchanges the $1$s in the permutation matrix at positions $(v(j),j)$ and $(i, v^{-1}(i))$ with $1$s at positions $(i,j)$ and $(v(j), v^{-1}(i))$. As a result, $\rk_{t_{v(j),i}v}(a,b) = \rk_v(a,b) +1$ if and only if $i \geq a$, $j \leq b$, $v(j)<a$ and $v^{-1}(i) >b$ (or in other words, if $(i,j) \in \SW(w)$ and . 
    the $1$s in that row and column lie outside of $\SW(w)$.
    Otherwise, the subrank stays the same. 

    \begin{enumerate}
        \item Assume $\rk_v(a,b) <m$. Then, for all $(i,j) \in \Dc(v)$, we see that $\rk_{t_{v(j),i}v}(a,b) \leq m$ and therefore $t_{v(j),i}v \leq w$. Hence, in this case, there are no unexpected zeros and the complexity of $\Nvw$ is given by $\cn(\Gvwt) - |\Dc(w)|$.
    \item Assume $\rk_v(a,b) =m$. Then, the only way to obtain $t_{v(j),i}v \not\leq w$ is if $(i,j) \in \SW(w)$, $v(j) <a$ and $v^{-1}(i) >b$. Since $\rk_v(a,b)= m$, there are $l$ rows $i \geq a$ with $v^{-1}(i) >b$ and $k$ columns $j \leq b$ with $v(j) <a$. These tuples $(i,j)$ must lie in the Rothe diagram and they are unexpected zeros. Hence, there are $lk = |\Dc(w)|$ unxepected zeros in $\Sigma_v$ and the complexity of $\Nvw$ is given by $\cn(\Gvwt)$.\qedhere
    \end{enumerate}\end{proof}

\end{document}
